\chapter{Amalie's Discrete Revolution}

\vspace{-\baselineskip}
\begin{minipage}[t]{\textwidth}
\lettrine{E}{mmi} Noether held a fascination for symmetry and conservation laws inherited from her illustrious namesake, Amalie Emmy Noether.  

\begin{figure}[H]
\includegraphics[width=1.0\linewidth]{Illustrations/vign-vign-emmi.png}
\caption{Amalie after her sister Emily (walking away) tells her she'd dumped Eminem}
\end{figure}
Her own pursuits had strayed from the Continuum, that had served her Amalie so well, to the digital, the graveyard of many an analyst, where the infinite continuity of classical physics rubs up (the wrong way) with the discrete, finite edges of integer lattices.
\end{minipage}

Amalie's program sought to reformulate reality on a grid, where finite differences and not partial derivatives, governed motion. Where rational numbers replace the infinitudes of analytic calculus. Her work began with the lattice of Gaussian integers, the set of complex numbers of the form \(a + bi\) (for integers \(a, b\)), a geometry of countable infinities. And she formed a coda:

\begin{quote}
If the discretizing program of physics arises from the need to regularise, to mathematically control the infinities of QED and QCD, and then to renormalise, redefining physical parameters to yield finite predictions, then space and time themselves may be better conceived as lattices. In that case, the arithmetic of Gaussian or Eisenstein integers may not merely model space-time but constitute its natural form, making regularisation redundant from the outset.
\end{quote}

Amalie imagined particles hopping between points on the Gaussian lattice, their paths governed by difference equations rather than differential ones. Energy and momentum conservation reformulated in terms of shifts operators that moved particles actively across the grid or passively by moving the grid itself. For Amalie, the classical tools of physics: calculus, the Lebesgue integral, and differential equations, were born from the convenient lies of the irrationals of analysis that are smoothness and continuity. But nature, doesn’t appeaer to submit to infinite divisibility. Amalie was intrigued by how power laws could be described in the discrete world. In calculus, a polynomial of degree \(n\) is governed by its \(n\)th derivative. In the lattice world, the \(n\)th difference of a sequence replaced derivatives:
\[
\Delta^n y_k = \text{constant}.
\]
This meant that the structure of the lattice itself encoded the laws of motion and change. The homogeneity of higher differences and how they remained invariant under shifts was analogous to the conservation laws of classical physics. Differences didn’t just describe the geometry of the lattice; they encoded its symmetry.

%%%%%%%%%%%%%%%%%%%%%%%%%%%%%
% >>>>>> ANNEX A (cut-sheet v2): Pascal's triangle and its Worpitsky extension (duplicate of the canonical Annexe A treatment) — block commented out below, pointer left in text <<<<<<
\textit{(Pascal's triangle and its Worpitsky extension (duplicate of the canonical Annexe A treatment) --- the full working is set out in Annexe A.)}

% \subsection*{Pascal’s Triangle and its Worpitsky Extension}
% 
% Amalie's thoughts turned to Pascal’s Triangle, a simple structure that embodies both the discrete and the infinite:
% \[
% \begin{array}{ccccccc}
%  & & & 1 & & & \\
%  & & 1 & & 1 & & \\
%  & 1 & & 2 & & 1 & \\
% 1 & & 3 & & 3 & & 1 \\
% \end{array}
% \]
% But what if, Amalie wondered, it could be extended to encode row numbers and sum products? She vaguely recalled an idea attributed to Worpitsky: each number in Pascal’s Triangle could be weighted by its row number, creating a richer lattice of sums and products.
% 
% \begin{figure}[H]
% \includegraphics[width=1.0\linewidth]{Illustrations/pyramid.png}
% \caption{All pyramids crumble in the end}
% \end{figure}
% 
% Finding such an extension of Pascal’s Triangle wasn’t just an exercise in Combinatorics. For Amalie, it represented a way to encode physical laws into the discrete framework of her focus. Each number in the triangle would become a building block for finite difference equations, while the row weights introduced an additional dimension of scaling and symmetry.
% 
% \lettrine{I}{t} was Amalie who first resurrected the Worpitsky Triangle into focus, where the triangle’s subscripted entries carried an added layer of organization and meaning, providing a framework for solving simultaneous linear equations and understanding polynomial sequences: 
% 
% \begin{itemize}
%     \item Start with the first row: \( 1_1 \).
%     \item The second row: \( 1_1 + 1_1 = 1\times1 +1\times1= 2_2 \).
%     \item The third row: \( 1_1 + 1_2 = 1\times 1+1\times 2=3_2 \).
%     \item The fourth row: \( 3_2 + 2_3 = 3\times 2+2\times 3= 12_3 \).
% \end{itemize}
% The entry \(W(k, j)\) in row \(k\) and column \(j\) of the triangle is determined by the sum-product rule applied to adjacent entries from the row above. The triangle up to the 5th row is thus:
% \[
% \begin{array}{ccccccccccc}
% & & & & & 1_1 & & & & &\\
% & & & & 1_1 & & 1_2 & & & &\\
% & & & 1_1 & & 3_2 & & 2_3 & & &\\
% & & 1_1 & & 7_2 & & 12_3 & & 6_4 & &\\
% & 1_1 & & 15_2 & & 50_3 & & 60_4 & & 24_5 &\\
% 1_1 & & 31_2 & & 180_3 & & 390_4 & & 360_5 & & 120_6 \\
% & & \vdots & & & \vdots & & & \vdots & & \\
% \end{array}
% \]
% 
% \noindent\
% Once you have the row \( W(k, j) \), the sum of \(k^{\text{th}}\) powers up to \(n\) is:
% 
% \[
% \sum_{i=1}^{n} i^k = \sum_{j=0}^{k} W(k, j) \cdot \binom{n}{j+1}.
% \]
% 
% \noindent\textbf{Example for \(k = 4\):} From the Worpitzky triangle:
% \[
% W(4, 0) = 1, \quad W(4,1) = 15, \quad W(4,2) = 50, \quad W(4,3) = 60, ...
% \]
% \begin{align*}
% \sum_{i=1}^8 i^4 &= 1\binom{8}{1} + 15\binom{8}{2} + 50\binom{8}{3} + 60\binom{8}{4} + 24\binom{8}{5} \\
% &= 8 + 420 + 1400 + 1575 + 672 = 4075.
% \end{align*}
% Analogous to Newton's interpolation on Pascal's triangle, Worpitzky provided an explicit formula for powers:
% 
% \[
% n^k = \sum_{j=0}^{k-1} W(k-1, j) \binom{n + j}{k}
% \]
% where \(W(k-1, j)\) are the Worpitzky numbers paralleling Newton's binomial-based forward difference expansion:
% 
% \[
% f(n) = \sum_{j=0}^k \Delta^j f(0) \binom{n}{j}
% \]
% but adapted specifically for polynomial sequences of powers, with \(W(k-1, j)\) replacing the role of finite differences, and with binomials shifted to \(\binom{n + j}{k}\).
% 
% \subsubsection*{Physical Reality on a Grid}
% 
% Amalie saw how finite differences could describe any polynomial law. For example, a quadratic relationship like \(y_k = k^2\) would yield a constant second difference:
% \[
% \Delta^2 y_k = 2.
% \]
% This made the infinite behavior of polynomials manageable within the finite structure of the lattice. She imagined discretized versions of physical laws: Maxwell’s, Einstein’s field equations in the spirit of Regge's Calculus. Each could be reformulated in terms of finite differences and lattice symmetries, removing the infinities of the continuum.
% 
% \smallskip
% 
% From the Gaussian integers to Pascal’s Triangle, Amalie saw a world where discrete structures could model the complexities of reality, constrained yet profound.
% 
% \begin{quote}
% "Mathematics," she thought, "must be rooted in the physical. And the physical, in turn, must encode mathematics."
% \end{quote}
% 
% >>>>>> end of annexed block <<<<<<
\section*{Emily being Discrete}

Eleanor leans back in her chair, studying Emily over the rim of her coffee cup. The café is murmuring, laden with sufficient incidentals to invite a confession or two. 

\smallskip

"So, Emily, what do you make of the investors circling us?" Eleanor asks, seemingly casual, but Emily knows better. Eleanor never wastes words. 

\smallskip

Emily swirls the last remnants of her espresso, eyes fixed on the liquid’s slow spiral. 

\smallskip

"Which ones specifically?" She asks, stalling.

\smallskip

"The ones with real money," Eleanor replies, arching an eyebrow. "The ones that want something in return. IP protections, conditions—more structure than we’d like." 

\smallskip

Emily forces a measured nod. She needs this investment more than she wants to admit. If the money is solid, if the funding is sufficient, it could relieve her of the impossible burden she’s carrying. The second mortgage, the quiet fear that Amalie sees too much, the slow realization that she’s playing a role she no longer believes in. 

\smallskip

"It’s a matter of trade-offs, isn’t it?" Emily says, choosing her words carefully. "If we accept conditional funding, it guarantees stability, but there has to be a cost? Losing some of our autonomy? I am someone who is pretty self-driven but I sometimes welcome an exogenous hit?" 

Eleanor nods, but her gaze sharpens. "I sense hesitation, Emily. This isn’t just an academic argument to you."

A pause. Emily swallows. The last thing she needs is scrutiny, the kind that might connect too many dots. 

\smallskip

"I do think we should know exactly what we’re getting into before we sign anything," Emily says smoothly. "Some of these investors will be looking for an exit strategy before we’ve even published. Others might push us into areas we might never want to explore. But very often the thrill really lies wandering down those dark alleys." 

\smallskip

Eleanor studies her, not even hiding it, then exhales. "That’s exactly what I wanted to hear. You see, Emily, I have my concerns too. I need people in the circle who can see beyond the mathematics, who understand the chase but are not naive to the nature of power, of influence."

\smallskip

Emily’s stomach knots. Influence. Power. Control. That is exactly what she is trying to regain in her own life. And yet, she nods. Smiles. Playing the role.

\smallskip

"Of course," she says. "We have to be careful. Always."

\smallskip

But inwardly, the word control echoes. She thinks of the investors as masters in their domain, yes, but only so because others submit. It’s recognition, yes, but filtered through dependence, through need. She wonders if that recognition ever satisfies. Probably not. It’s hollow if it comes from someone whose autonomy is compromised\footnote{Hegel’s dialectic of Master and Slave suggests that the Master seeks recognition from the Slave, yet the recognition is incomplete because it comes from someone not fully free. See G. W. F. Hegel, Phenomenology of Spirit, trans. A. V. Miller (Oxford: Oxford University Press, 1977), §189.}.

\smallskip

She swirls the cold dregs of espresso again, watching the spiral tighten, thinking: who’s the master here, really? 

"Of course," she says. "We have to be careful. Always."

And for now, Eleanor believes her because she is doing a good job of pretending to herself of feeling so ambivalent. But how long until she doesn’t?

\subsection*{Emily Contemplates the Circle as a Fractured Whole}

Eleanor pours them both another drink and takes a temporary leave such is the will of the body within. Emily now alone, overriding the appeal of her collective nerve-endings, resists the urgent, impulse to reach for the phone: 
that outsourced memory, that copter of her own circling thoughts. She has to resist hard for the collective within is strong, she knows better than that mindless crowd that less is more. The dimming reserve of the cafe, the murmur of the city networking outside interlacing the static charging her mind. She, her attended mind is just thinking.

\smallskip


She has spent years navigating the precarious balance between mathematical purity and financial pragmatism, between the pursuit of truth and the necessity of survival. But now, after what has happened, after what she has concealed, \textit{what even is she}?

\smallskip

The Circle, a group bound by a shared kind of intellect, should be something greater than its individual members. And yet, standing now on the periphery, Emily sees something in herself fracturing from within:

\begin{quote}
\textit{no longer is she of coherent mind, merely a tension of wills, an organism fighting against itself}. 
\end {quote}

\smallskip

A body is made up of cells, each specialized, each following its own local rules, yet somehow, against all statistical odds, \textit{an individual emerges}. A self, a mind, a will. But what happens when the cells disagree? What happens when one cell prioritizes its own survival over the integrity of the whole? Cancer. The rogue cell, multiplying unchecked, begetting offspring: the ultimate irony is their seeking immortality save for the demise of the host organism their proliferation imperils.

\smallskip

\textit{And is she that rogue cell?} She had told herself it was survival. That she had no choice but to stay silent about the coming collapse of the Invariant Fund. That her \textit{fiduciary duty}, (really what does that even mean?) her position, her \textit{clients}—all of these had to come before the Circle, before her own sister.

\smallskip

\textit{But was it really necessity? Or was it self-preservation?}

With the quiet echo of her \href{https://kingsoakacademy.clf.uk/our-academy/kingswood-old-scholars-association/}{school motto}: \textit{esse non videri}, to be, not to be seen, still etched in her mind, it had become a cipher for a deeper paradox. We are chemical soups, endochrinal symphonies, orchestrated not by a single conductor but by a multitude of semi-autonomous instruments. Our glands whisper instructions into our bloodstreams, invisible and persistent, shaping desire, fear, hunger, and memory. These signals don’t ask permission; they just fire, they compel, they fade. And now, with tensor silicon chips (human's synthetic neurons) we attempt to simulate this distributed intelligence. To clone the subtlety of the biochemical ballet that is us in lines of machine code on tapestries of transistors. We chase the illusion of selfhood in our creations, believing intelligence arises from centrality. Not likely.

\smallskip

From the bottom up, we are assemblies. Not one voice but many, coordinated just enough to give the impression of a single “I.” From the top down, we impose stories of unity: of agency, of will, of continuity. Our nervous systems are reconciliation engines, stitching coherence from contradiction through feedback, inhibition, emergence. Like a hive, like a flock, like a city at dawn, we operate through distributed cognition. Emily wasn’t just a self: she was selves, harmonized. To be, not to seem, was a myth too. Even being is a seeming, if you dare to look closely enough.

\smallskip

Perhaps the convenient myth we tell each other is that a collection of functionally distinct, purposeless parts can coalesce into a unified, purposeful whole. Somewhere just beyond the sophistication of slime mould, something emergent stirs: coordination mistaken for intention, pattern mistaken for plan. That beyond chemical whispers and synaptic noise, meaning arises, not because it’s there, but because we are built to see it. This is the Bruce Hoodian \textit{delusion} of selfhood: from the bottom up, we are distributed systems of semi-autonomous components, yet top down, we act as if we are singular beings. Our nervous systems reconcile these contradictions through hierarchical coordination, feedback loops, and constraints. That is, just as a society, a hive, or even a flock of birds does.

\smallskip

Yet within this unity, tension is always latent. Sometimes a rogue cell, like cancer, disrupts the delicate balance. Sometimes an individual resists the will of the collective —be that collective a family, a research group, or an ideal.

\smallskip

And yet—Emily wonders—not all resistance is malignant. But what if resistance calcifies? What if, in rationalizing personal necessity, she had become utilitarianism’s shadow—its monster? The cold calculus that demands sacrifice, the equation that never stops balancing itself, until even the self is expendable for the 'greater good.' She recalls Parfit’s warning: the Utility Monster, the being who can convert resources into happiness so efficiently that, by utilitarian logic, it justifies absorbing everything and everyone else\footnote{Derek Parfit introduces the concept of the Utility Monster in Reasons and Persons (Oxford: Oxford University Press, 1984), p. 336, to critique utilitarianism's potential to justify extreme inequality if it maximizes overall utility.}.

\smallskip

What if the Invariant Fund with its roots in the Circle, in its collective genius, was itself becoming that monster? Hungry, optimizing, consuming any dissent, any weakness, any individuality in the name of a collective advance. And what if, by remaining silent, she wasn’t just a rogue cell but fuel?

\smallskip

She presses her temples lightly, as if that might soothe the fracture she feels widening within. A tension of wills. Not coherent, not unified. Not even truly divided. Just... multiplying.

\textit{So, how do you remain “you” when your parts are in such flux?} Perhaps identity itself is no singular flame but a quorum of impulses and reconciliations, bickering towards coherence.

\smallskip

She thinks of the Circle, not as a club, not a cabal, but a quiet \textit{quorum} of minds, tending to the mathematics that becomes machine code, the assembly that speaks to firmware etched into the substrate of the lattice of optical fibres and servers that cradle the modern world. She alone cannot be entrusted to toy with such power. No one can. But perhaps a chorus of flawed, watchful minds: a fragile murmering of like minds, can be. And maybe the Circle, like the body, fails not from dissent, but from losing the capacity to resolve its inner differences. That is, like a nervous system gone static, no longer ecstatic, unable to tell signal from noise. A thought that hangs over her like a theorem she’s too tired to prove.

\smallskip

% ============ COMMENTED OUT (cut-sheet v2): paragraph duplicated in 9_7meetings.tex ('The Circle—mathematicians, logicians…') ============
% The Circle of mathematicians, logicians, theorists, each a node in an intricate web of thought, has never been about any singular ambition. For Eleanor, it is a bastion of uncorrupted intellectualism. For Willard, a monument to pure mathematics, untouched by vulgar concerns of money. For others, it is a means to personal relevance, to securing the resources that keep their research afloat. They do not have a single purpose. They never did.
% ============ end commented block ============

\smallskip

A society is no different. It appears to function as a whole, but only because its competing forces reach a temporary balance. The moment that equilibrium shifts, everything unravels. Emily had thought of the Circle as something greater than its members; something that transcended individual ambitions.

But what if that was the illusion? A human being is little more than the sum of its performative organs: it is the silent compositional harmony of those organs working together. Yet the same system that holds the body together can betray it. A seizure. A misfire. A catastrophic breakdown in coordination.

Is that what is happening now? Is the Circle experiencing its own kind of neural collapse?

She recalls a fragment of Heidegger, or maybe just the impression of it: that humans are always seeking a purposeless purpose, a striving without a clear telos, just to keep being\footnote{Heidegger explores this in Being and Time, where Dasein is characterized by an ever-unfolding projection of possibilities—Sein-zum-Ende, being-towards-death—without ultimate closure. See Martin Heidegger, Being and Time, trans. John Macquarrie and Edward Robinson (Harper \& Row, 1962), p. 294.}. This drive is not imposed from above but emerges from the bottom up just like a society whose purpose is not ordained but arises simply because coordinated survival requires it. But when coordination fails, when some parts defect or malfunction, the system turns on itself.

Maybe that’s what the Circle is now. A seizure in slow motion. An organism convulsing on the edge of its own limits, generating movements that look like intention but are, in fact, just spasms. And if Emily is honest, she no longer knows whether she is the signal or the noise.

She thinks of her sister, of Eleanor’s tireless efforts to preserve the Circle’s ideals. Does Eleanor even know? Would she look at Emily now and see a friend, a sister, or a traitor? She exhales sharply pondering the inconvenience of her \textit{unseen being}:

\begin{quote} We long to believe we belong to a greater collective from without, yet here we are always living out the consequences of biochemical decisions made within. We are mere stereochemistry, writ (slightly) complex. \end{quote}

Surely the Circle will endure ieven if the fund doesn't in some form, but Emily wonders what will be left of herself inside it. 


% ============ COMMENTED OUT (cut-sheet v2): verbatim duplicate of 'Emily's Realization' — kept version lives in 9_7meetings.tex ============
% \subsection*{Emily’s Realization: The Circle Was Never Pure}
% 
% As Emily returns home she catches herself leaning against the heavy oak table, fingers pressed into the grain, trying to anchor herself to something real. She had always known the Circle was a loose network—a construct of alliances and ambitions. But she had convinced herself it had true purpose—something more ideal than even her own idealism.
% 
% She sees now that she had been kidding herself.
% 
% \smallskip
% 
% The others had already formed a thesis of pragmatism. Somewhere, away from Eleanor’s principles, the group had quietly acknowledged the unspoken truth: mathematics alone would never be enough.They had acted on that knowledge by creating an investment fund with purpose—an attempt to wield their insights upon the world of finance. Not merely for wealth, but to demonstrate: that the abstract could move the real.
% 
% \smallskip
% 
% And yet, it had failed. Chris and Willard’s theories of order within chaos had cracked. The Invariant Arbitrage system, for all its mathematical elegance, collapsed under the weight of the uncontrollable. Chris’s models treated price movements as lattice structures — probabilistic near-certainties woven into a seemingly inevitable profit engine. But the markets weren’t a purely endogenous emergent order, waiting to be reverse-engineered in real time. They were shaped—intervened upon—by exogenous forces. Social. Political. Intentional.
% 
% \smallskip
% 
% Bond vigilantes. Investment banks. Central banks and their unspoken alignments behind closed doors.
% 
% \subsection*{The Unresolvable Paradox: Bottom-Up vs. Top-Down}
% 
% The failure of the fund had exposed their naïve assumption: that markets, like physical systems, would resolve into knowable structures. But finance was not governed by natural law. It was not an emergent phenomenon of impartial agents—it was a game, and those with power wrote the rules.
% 
% \smallskip
% 
% If the Circle had truly been a sanctuary for pure mathematics, the fund would never have existed. The moment they sought to monetize mathematics, they forfeited their autonomy. They had already bent the knee to pragmatism.
% 
% \smallskip
% 
% Eleanor, with her unwavering belief in the purity of research, had remained blissfully unaware. Or perhaps they had only convinced themselves she was. How much longer could that last? Would she still believe in the collective dream, knowing it had already splintered?
% 
% \smallskip
% 
% Emily exhales, trying to silence the churn of her (owned) mind. She sees the contradiction clearly now:
% 
% \begin{quote} If she tells Eleanor, the Circle shatters.
% If she stays silent, it rots from within. \end{quote}
% 
% But the rot, she realizes, had begun long ago—the moment they dabbled in finance. Maybe Eleanor had always known. Maybe they all had, and chose not to see it.
% 
% For now, Emily will say nothing. Not out of cowardice, no, but because—for the first time—she’s no longer certain who she is loyal to.
% 
% 
% 
% ============ end commented block ============



